\section{Discussion}
\label{sec:discussion}

%\subsection{Further strategies}
%\todo{More strategies?..}
%\subsubsection{Multiple Type Strategies}
%\todo{FloatAndInt}
%\subsubsection{Strategies for Dictionaries}
%\subsubsection{Run-Length Strategies}
%\todo{Ranges, length-of-nil, ...?  We don't have this, so future work...}

The collection types supported by rstrategies are fixed sized and variable sized lists.
Of course there are may more collection types and algorithms used in modern programs.
The collection framework in the Java standard library features 8 interfaces and 10 implementation classes.
A modern Squeak image \footnote{\url{http://www.squeak.org/}} has as many as 79 subclasses of \lstinline{Collection}, because it includes collections with very special purposes.
However, rstrategies is a library for optimizing collections on VM level, which means it is only meaningful to support collections that occur as \emph{primitive types} in multiple programming languages.
Even though Java supports collections like \lstinline{ArrayList}, \lstinline{TreeSet} and \lstinline{HashMap}, the only primitive collection type is the \lstinline{Array}.
The same goes for Smalltalk and many other programming languages.

It would be possible to take the approach of storage strategies even further and let the VM transparently optimize collection types from the standard library of the programming language.
For example, since the interface and semantics of Javas \lstinline{ArrayList} are well defined, the VM could try to handle instances of \lstinline{ArrayList} by itself, without actually executing the Java bytecode in the methods of the \lstinline{ArrayList} class.
This might lead to improved performance when accessing elements of \lstinline{ArrayList}, but this approach has several drawbacks.
One important function of a standard library is to reduce the complexity of the VM.
Moving parts of the library into the VM would only further increase its complexity.
Another purpose of having a separate standard library is that it is able to evolve independently of the VM.
It is very hard to exactly mirror the code of the standard library and the suggested approach would require releasing a new VM version after every substential change to the standard library.

Therefore we chose to only optimize primitive collection types with rstrategies.
The rest of this section lists and discusses several collection types not supported by rstrategies, although they do occur as primitive types in several modern programming languages.

\subsection{Non-List Collections}

We examined the primitive datatypes of the top programming languages according to the TIOBE Index for July 2014 \footnote{\url{http://www.tiobe.com/index.php/content/paperinfo/tpci/index.html}}.
Even though such an index is never exact and changes fairly quickly, we believe it provides a satisfactory view of the state-of-the-art of programming languages.
We looked at the top 10 programming languages which are usually executed by a VM.
C and C++ are excluded, for example, since storage strategies are not applicable for statically compiled languages.
Table \ref{tab:collectionTypes} summarizes the results by listing the primitive datatypes of each language and mapping them to abstract collection types which are explained in the following.

\begin{table}
\centering

\begin{tabular}[ht]{ l l l l }
\toprule
Language & Fixed List & Variable List & Dictionary \\
\midrule
Java & Object, Array & {} & {} \\
Objective-C & Object, Array & {} & {} \\
C\# & struct, Object, Array & {} & {} \\
PHP & {} & array & array \\
Python & tuple, bytearray & list & object, dict \\
JavaScript & {} & object & object \\
Perl & Object, array & {} & hash \\
Visual Basic .NET & Object, Array & {} & {} \\
F\# & list, tuple, array & {} & struct, record \\
Ruby & {} & array & object, hash \\
\bottomrule
\end{tabular}

\caption{Programming languages and their primitive collection types.}
\label{tab:collectionTypes}
\end{table}

The first line of table \ref{tab:collectionTypes} lists 3 abstract collection types: \emph{Fixed List}, \emph{Variable List} and \emph{Dictionary}.
\emph{Fixed List} and \emph{Variable List} correspond to the collection types supported by rstrategies, using the mixin hierarchy described in \ref{sec:solutionOptimizations}.
Therefore, all collection types listed in these two columns could be implemented and optimized using rstrategies.

\emph{Dictionary} is a third abstract collection type that is quite important for many languages.
There are two kinds of \emph{Dictionary} based collections.
The first kind corresponds to Pythons or Rubys \lstinline{object} type, which only allows symbols as keys.
This type is an integral part of these languages and therefore its optimization is central to VM performance.
The standard optimization for such \emph{slot objects} is called maps, has been first developed for the Self VM \cite{self, selfDissertation}.
This advanced optimization is not part of rstrategies because it is inherently language dependent.
The second kind of \emph{Dictionary} based collections are regular lookup tables as implemented in standard libraries of many languages.
This collection type is supported as built in type in Python, Perl and Ruby.
Other languages, like JavaScript, feature types which can be used as both slot objects or general lookup tables.

Storage strategies can be extended to include support for lookup tables.
This has been done for the Python interpreter PyPy \cite{strategies} and in the JRuby+Truffle project \cite{smallDataStructures}.
We have not yet implemented lookup tables as part of rstrategies yet, but intend to do so as part of our future work on this library.

The contents of table \ref{tab:collectionTypes} are complete for every language, with one exception.
Python has several additional primitive types for special use cases: \lstinline{buffer}, \lstinline{xrange}, \lstinline{set}, \lstinline{frozenset}, \lstinline{iterator}, \lstinline{generator}.
This variety of primitive collection types is special to Python and all these types can be mapped to the \emph{Fixed List} or \emph{Variable List} types by adding additional algorithms to the \lstinline{store} and \lstinline{fetch} routines.
Therefore we chose not to support special datatypes like these.


\begin{comment}
% This "idea" is not really meaningful
\subsection{Multilevel-Strategies}

In theory, storage strategies could even be applied to multiple levels, like illustrated in figure \ref{fig:MultipleLevels}.
In this example, there are multiple equally sized lists of integer objects, which in turn are collected in a single, larger list.
This two-level setup could be optimized by placing the unboxed version of all integers in consecutive memory.

We decided to not take the storage strategies approach that far.
For one thing, scenarios like that are rare.
Furthermore, this optimization would require the VM to store and constantly evaluate additional control information, like the number of levels optimized and the length of the contained lists.
This requires additional resources and the required analysis would likely be too expensive to pay off.

\begin{figure}[ht]
\centering
\includegraphics[width=0.7\textwidth]{strategies_multipleLevels}
\caption{In theory, even multiple levels of storage can be optimized.}
\label{fig:MultipleLevels}
\end{figure}
\end{comment}

\subsection{Non-Contiguous Lists}

Functional programming languages often feature a primitive datatype for linked lists.
Lisp for example has the built in type \lstinline{cons} which represents one link of a linked list, formed by two fields: the value (named \lstinline{car}) and the pointer to the next link (named \lstinline{cdr}).
While these labels have historical reasons originating in the 1950s, the linked list is still a central data structure in modern Lisp programs.
Since accessing arbitrary elements of a linked list requires walking the entire list up to that element, optimizing linked lists would be very beneficial.
For example, all elements of the list could be placed in consecutive memory, allowing constant access to any element.
However, such optimizations bring about further complications and are a separate (although related) research topic to storage strategies.
Therefore, we consider optimizations of linked lists out of scope for rstrategies.
